\documentclass[11pt]{article}
\usepackage[margin=1in]{geometry}
\usepackage{amsmath,amssymb,amsthm,mathtools}
\usepackage{graphicx}
\usepackage{microtype}
\usepackage[hidelinks]{hyperref}

\newtheorem{theorem}{Theorem}
\newtheorem{lemma}[theorem]{Lemma}
\newtheorem{remark}[theorem]{Remark}
\newcommand{\R}{\mathbb R}
\newcommand{\N}{\mathbb N}
\newcommand{\cG}{\mathcal G}
\newcommand{\esssup}{\operatorname*{ess\,sup}}
\newcommand{\dens}{\operatorname{dens}}

\title{Dense Yes, Closed No for Algebras of Nowhere-Bounded Functions}
\author{Candidate full solution packet\\\small arXiv:1207.3818, Problem 3}
\date{11 August 2026}

\begin{document}
\maketitle

\noindent\textbf{Status.} Candidate full answer in the source's intended
cardinal/separable setting, pending expert review: $\cG_p\cup\{0\}$ contains
a dense free algebra, but it contains no nonzero $L^p$-closed subalgebra.  The
closed-algebra obstruction is valid on every measure space and for every
$0<p<\infty$.

\section{Source problem and setting}

G\l ab, Kaufmann, and Pellegrini define
\[
 \cG_p=\{f\in L^p(\mu):f\text{ is nowhere }L^\infty(\mu)\}
\]
and ask the following on printed page 17 of \cite{GKP}:

\begin{center}
\includegraphics[width=0.98\textwidth]{figures/open_problem_crop.png}
\end{center}

The positive part below uses exactly the factorial dense partition constructed
in the proof of Theorem 1.7(e) of \cite{GKP}.  Namely, there are pairwise
disjoint Borel sets $(B_j)_{j\geq1}$ such that
\begin{equation}
  \mu(B_j)\leq\frac1{j!},
  \qquad
  \mu(U\cap B_j)>0
\end{equation}
for every $j$ and every nonempty open $U\subset X$.  The source obtains these
sets when $X$ has a countable $\pi$-base and $\mu$ is atomless, outer regular,
and has full support.

For density, we make explicit the same cardinal convention used in the
source's dense-lineability proof:
\begin{equation}
  \dens L^p(\mu)\leq\mathfrak c.
\end{equation}
This holds in particular whenever $L^p(\mu)$ is separable, including
$[0,1]$ with Lebesgue measure and the standard examples highlighted in the
paper.  The source itself invokes separability/cardinality $\mathfrak c$ in
its proof of Theorem 1.7(d).

\section{Proof intuition}

The source already provides free generators
\[
  g_\theta=\sum_{j\geq1}\theta^j\chi_{B_j},\qquad \theta>1.
\]
Factorial decay makes every polynomial in them $p$-integrable, while each
open set sees all the exponentially growing levels.  Their algebra is not
dense because these functions are constant on every $B_j$.  To repair this,
add bounded $L^p$ functions chosen from a dense family and make the
pathological perturbations arbitrarily small in $L^p$.  Rational independence
of the logarithms of the bases gives every nonzero polynomial a unique largest
exponential rate.  Bounded perturbations only create smaller rates, so this
leading term cannot cancel on $B_j$.

The negative direction is more general.  In a closed subalgebra $A\subset
L^p$, fix $f\in A$.  A Baire argument shows that multiplication by $f$ is a
bounded operator on $A$, even for $p<1$.  Applying its norm estimate to the
powers $f^n$ gives a uniform bound on $\|f\|_{np}$ as $n\to\infty$, which
forces $f\in L^\infty$.  Thus every closed $L^p$ subalgebra consists entirely
of bounded functions and cannot have a nonzero member of $\cG_p$.

\section{The dense free algebra}

\begin{theorem}\label{thm:dense}
Let $0<p<\infty$.  Assume (1) and (2).  Then
$\cG_p\cup\{0\}$ contains a dense free real algebra in $L^p(\mu)$.
\end{theorem}

\begin{proof}
Bounded $L^p$ functions are dense in $L^p$: truncate any representative in
absolute value.  Choose a dense family
\[
  D\subset L^p(\mu)\cap L^\infty(\mu),
  \qquad |D|\leq\mathfrak c.
\]
Index a family by $I=D\times\N$.  Since $|I|\leq\mathfrak c$, choose numbers
$(\theta_i)_{i\in I}\subset(1,\infty)$ such that
$(\log\theta_i)_{i\in I}$ is linearly independent over $\mathbb Q$.  Put
\[
  g_i=\sum_{j=1}^\infty\theta_i^j\chi_{B_j}.
\]
By (1),
\[
 \|g_i\|_p^p
 \leq\sum_{j=1}^\infty\frac{\theta_i^{pj}}{j!}<\infty.
\]

Write $i=(h,n)$ with $h\in D$ and $n\in\N$, choose $\varepsilon_i>0$ so
small that
\[
  \|\varepsilon_i g_i\|_p<\frac1n,
\]
and define
\[
  z_i=h+\varepsilon_i g_i.
\]
For every fixed $h\in D$, the sequence $(z_{(h,n)})_n$ converges to $h$ in
$L^p$.  Hence the set of generators $(z_i)_{i\in I}$, and therefore the
algebra it generates, is dense in $L^p$.

We prove that every nonzero polynomial without constant term in finitely many
distinct generators belongs to $\cG_p$.  Let
\[
 P(t_1,\ldots,t_r)=\sum_{\alpha\in F}c_\alpha t^\alpha,
 \qquad
 F\subset\N_0^r\setminus\{0\},
\]
where all $c_\alpha\ne0$.  For the selected generators set
\[
 \lambda_\alpha=\prod_{\ell=1}^r\theta_{i_\ell}^{\alpha_\ell}.
\]
Rational independence of the logarithms makes the $\lambda_\alpha$ pairwise
distinct.  Let $\alpha^*$ be the unique index with largest $\lambda_\alpha$.

On $B_j$ we have
\[
 z_{i_\ell}(x)=h_\ell(x)+\varepsilon_\ell\theta_{i_\ell}^j,
 \qquad \|h_\ell\|_\infty<\infty.
\]
Expanding the polynomial, its unique term with the largest exponential rate
is
\[
 c_{\alpha^*}\varepsilon^{\alpha^*}
 \lambda_{\alpha^*}^j,
 \qquad
 \varepsilon^{\alpha^*}:=
 \prod_{\ell=1}^r\varepsilon_\ell^{\alpha^*_\ell}.
\]
Every other expanded term has a rate strictly smaller than
$\lambda_{\alpha^*}$ and a coefficient uniformly bounded in $x$.  Since only
finitely many terms occur, there are $C>0$ and
$1\leq\rho<\lambda_{\alpha^*}$ such that, uniformly for $x\in B_j$,
\begin{equation}
 \left|P(z_{i_1},\ldots,z_{i_r})(x)
 -c_{\alpha^*}\varepsilon^{\alpha^*}
  \lambda_{\alpha^*}^j\right|
 \leq C\rho^j.
\end{equation}
It follows that for all sufficiently large $j$,
\begin{equation}
 |P(z_{i_1},\ldots,z_{i_r})(x)|
 \geq\frac12|c_{\alpha^*}|\varepsilon^{\alpha^*}
       \lambda_{\alpha^*}^j
 \quad (x\in B_j).
\end{equation}

For completeness, the polynomial is in $L^p$.  On $B_j$ its absolute value
is at most $C'(1+\Lambda^j)$ for fixed $C',\Lambda$, so (1) and factorial
decay make the $p$th-power integral over $\bigcup_jB_j$ finite.  Off this
union all $g_i$ vanish.  There the polynomial is a sum of positive-degree
monomials in bounded $L^p$ functions, and every such monomial is in $L^p$.
This argument does not use convexity and is valid for $p<1$ as well.

Finally, every nonempty open $U$ meets every $B_j$ in positive measure.  The
lower bound (4) tends to infinity, so the polynomial is not essentially
bounded on $U$.  It is therefore in $\cG_p$.  In particular no nonzero
polynomial relation exists, proving that the generators are free.
\end{proof}

\section{Why a closed algebra is impossible}

\begin{theorem}\label{thm:closed}
Let $(X,\Sigma,\mu)$ be any measure space and $0<p<\infty$.  If $A$ is a
pointwise subalgebra of $L^p(\mu)$ which is closed in the $L^p$ topology, then
\[
  A\subset L^\infty(\mu).
\]
Consequently $\cG_p\cup\{0\}$ contains no nonzero closed subalgebra.
\end{theorem}

\begin{proof}
Fix $f\in A$.  For $m\in\N$, let
\[
  E_m=\{g\in A:\|fg\|_p\leq m\}.
\]
Each $E_m$ is closed in $A$.  Indeed, if $g_k\to g$ in $L^p$, select a
subsequence converging almost everywhere and use Fatou's lemma to obtain
\[
  \|fg\|_p^p\leq\liminf_k\|fg_k\|_p^p.
\]
Moreover $A=\bigcup_mE_m$, because $A$ is an algebra and hence $fg\in A
\subset L^p$ for every $g\in A$.

The closed subspace $A$ is a complete metrizable topological vector space,
also when $p<1$.  By Baire's theorem, some $E_m$ has nonempty interior.
Thus there are $g_0\in A$ and $r>0$ such that
\[
  g_0+\{h\in A:\|h\|_p<r\}\subset E_m.
\]
Since also $g_0\in E_m$, the triangle inequality when $p\geq1$, and the
$p$-triangle inequality when $p<1$, give a constant $K_p$ such that
\[
  \|fh\|_p\leq K_pm\qquad(\|h\|_p<r).
\]
Scaling now yields a constant $C_f<\infty$ with
\begin{equation}
  \|fh\|_p\leq C_f\|h\|_p\qquad(h\in A).
\end{equation}

Apply (5) repeatedly to $h=f^{n-1}$.  Since all powers belong to $A$,
\[
 \|f^n\|_p\leq C_f^{n-1}\|f\|_p,
 \qquad\text{so}\qquad
 \|f\|_{np}\leq
 C_f^{(n-1)/n}\|f\|_p^{1/n}.
\]
We claim that $\|f\|_\infty\leq C_f$.  Otherwise choose $t>C_f$ such that
$E=\{|f|>t\}$ has positive measure.  Since $f\in L^p$, this measure is also
finite.  Therefore
\[
  \|f\|_{np}\geq t\,\mu(E)^{1/(np)}\longrightarrow t,
\]
contradicting (6).  Thus $f\in L^\infty$.  As $f\in A$ was arbitrary,
$A\subset L^\infty$.

If in addition $A\subset\cG_p\cup\{0\}$, any nonzero $f\in A$ would be
globally essentially bounded and nowhere essentially bounded simultaneously,
an impossibility.  Hence $A=\{0\}$.
\end{proof}

\section{Conclusion, scope, and novelty}

Combining Theorems~\ref{thm:dense} and \ref{thm:closed} gives the sharp
dichotomy requested in Problem 3:
\[
 \boxed{\text{dense free algebra: yes; nonzero closed algebra: no.}}
\]
The dense statement applies under the factorial partition used by the source
and the cardinal condition (2), in particular throughout the separable
setting explicitly used in the paper and for $L^p[0,1]$.  If the source's
displayed topological hypotheses are interpreted to include exotic $L^p$
spaces of density strictly above $\mathfrak c$, no claim about dense algebras
in that additional cardinal regime is made.  The negative closed-algebra
theorem is unrestricted.

The proof was audited separately at the unique-leading-rate step, the
off-partition integrability step, the $p<1$ Baire estimate, and the passage
from power bounds to essential boundedness.  The accompanying
\texttt{verification.md} records these checks; the finite script is only a
sanity check and is not used as proof.

A bounded novelty search on 11 August 2026 covered the four cheap run indexes,
the locally parsed arXiv corpus, and current arXiv-domain searches for the
exact title, arXiv id, Problem 3 wording, ``nowhere essentially bounded'',
``dense algebra'', and ``closed subalgebra''.  It recovered the source and its
predecessor arXiv:1204.6404, which state the question as open, but no later
answer.  Novelty remains subject to specialist review.

\medskip
\noindent\textbf{Human-review recommendation.} Review as a candidate full
answer in the paper's intended setting.  Check especially that condition (2)
matches the source's standing cardinal convention; independently, Theorem
\ref{thm:closed} is a fully general negative theorem.

\begin{thebibliography}{9}
\bibitem{GKP}
S. G\l ab, P. L. Kaufmann, and L. Pellegrini,
\emph{Large structures made of nowhere $L^p$ functions},
Studia Math. 216 (2013), no. 1, 51--69; arXiv:1207.3818,
Problem 3 on printed page 17.
\end{thebibliography}

\end{document}
